% !TEX program = xelatex
% 공통수학2 · Ⅰ. 도형의 방정식 · 3. 도형의 이동 · 01 평행이동 · 1/1차시
\documentclass[11pt,a4paper]{article}
\usepackage{kotex}
\usepackage{iftex}
\ifPDFTeX\else
  \setmainhangulfont{Noto Serif CJK KR}
  \setsanshangulfont{Noto Sans CJK KR}
\fi
\usepackage[margin=2.1cm]{geometry}
\usepackage{parskip}
\usepackage{enumitem}
\usepackage{array}
\usepackage{tabularx}
\usepackage{booktabs}
\usepackage{xcolor}
\usepackage{amssymb}
\usepackage{tikz}
\usepackage[most]{tcolorbox}
\usepackage{fancyhdr}
\usepackage{titlesec}

\definecolor{goldc}{HTML}{B07C22}
\definecolor{goldstrong}{HTML}{8A5F16}
\definecolor{indigoc}{HTML}{2B3B57}
\definecolor{indigosoft}{HTML}{6E82A6}
\definecolor{linec}{HTML}{D6DACB}
\definecolor{surface2}{HTML}{E4E8DC}
\definecolor{notebg}{HTML}{FBF3E3}
\definecolor{inksoft}{HTML}{52604F}

\titleformat{\section}{\normalfont\sffamily\bfseries\large\color{indigoc}}{\thesection}{0.6em}{}
\titlespacing{\section}{0pt}{16pt}{8pt}
\newtcolorbox{ideabox}{enhanced, breakable, colback=white, colframe=goldc, boxrule=0.5pt, leftrule=3pt, arc=2pt, left=10pt, right=10pt, top=8pt, bottom=8pt}
\newtcolorbox{calloutbox}{enhanced, breakable, colback=white, colframe=indigosoft, boxrule=0.5pt, leftrule=3pt, arc=2pt, left=10pt, right=10pt, top=7pt, bottom=7pt}
\newtcolorbox{notebox}{enhanced, breakable, colback=notebg, colframe=goldc, boxrule=0.5pt, leftrule=3pt, arc=2pt, left=10pt, right=10pt, top=7pt, bottom=7pt}
\newtcolorbox{answerbox}{enhanced, breakable, colback=white, colframe=linec, boxrule=0.5pt, arc=2pt, left=10pt, right=10pt, top=7pt, bottom=7pt}
\newcommand{\lessonbanner}[3]{%
  \begin{tcolorbox}[enhanced, colback=indigoc, colframe=indigoc, boxrule=0pt, arc=0pt,
    left=14pt, right=14pt, top=14pt, bottom=10pt, borderline south={2.4pt}{0pt}{goldc}, fontupper=\color{white}]
    {\sffamily\footnotesize\color{white!85!black} #1}\\[4pt]{\sffamily\bfseries\Large #2}\\[3pt]{\sffamily\small\color{white!88!black} #3}
  \end{tcolorbox}}
\pagestyle{fancy}\fancyhf{}\renewcommand{\headrulewidth}{0pt}
\fancyfoot[L]{\footnotesize\color{inksoft} 공통수학2 · 2026학년도 2학기 · 지평선고등학교}
\fancyfoot[R]{\footnotesize\color{inksoft} 교사용 지도안 -- 배포 금지}

\begin{document}
\lessonbanner{공통수학2 · Ⅰ. 도형의 방정식 · 3. 도형의 이동}{01 평행이동 --- 1/1차시}{점과 도형의 평행이동 · 교사용 지도안 (2026학년도 2학기)}
\vspace{10pt}

\noindent\begin{tabularx}{\linewidth}{@{}>{\bfseries\color{indigoc}}p{2.6cm}X@{}}
\toprule
대상 & 고1 · 2개 반 · 반당 13명 \\
차시 & 50분 (도입 10 · 전개 30 · 정리 10) \\
성취기준 & 평행이동을 탐구하고, 실생활과 연결하여 문제를 해결할 수 있다. \\
선행 학습 & 2중단원 --- 원의 방정식 \\
\bottomrule
\end{tabularx}

\section{학습 목표 \& Big Idea}
\begin{ideabox}
\textbf{\color{goldstrong}영속적 이해 (Big Idea)}\\
도형을 옮기면 그 방정식도 함께 `옮겨진다' --- 그런데 점의 좌표는 $+$로 이동하지만, 방정식은 반대로 $-$를 대입해야 한다는 미묘한 비대칭을 이해하는 것이 핵심이다.
\vspace{4pt}
\small\color{inksoft} 핵심개념: \textbf{이동} \quad\textbullet\quad 핵심개념: \textbf{불변량(모양)} \quad\textbullet\quad 일반화: \textbf{좌표의 변환과 방정식의 변환은 서로 반대 방향이다}
\end{ideabox}
\begin{calloutbox}
\textbf{차시 학습 목표}
\begin{itemize}[leftmargin=1.4em, itemsep=2pt, topsep=4pt]
  \item 점을 $x$축·$y$축 방향으로 평행이동한 점의 좌표를 구할 수 있다.
  \item 방정식 $f(x,y)=0$이 나타내는 도형을 평행이동한 도형의 방정식 $f(x-a,y-b)=0$을 구할 수 있다.
\end{itemize}
\end{calloutbox}

\section{질문의 연결}
\begin{tabularx}{\linewidth}{@{}>{\bfseries\color{indigoc}\small}p{2.4cm}X@{}}
사실적 질문 & 점 $(x,y)$를 $x$축 방향으로 $a$, $y$축 방향으로 $b$만큼 평행이동하면 좌표는 어떻게 바뀔까? \\[6pt]
개념적 질문 & 점은 $(x+a, y+b)$로 `더하는데', 왜 도형의 방정식은 $x$ 대신 $x-a$를 `빼서' 대입해야 할까? \\[6pt]
논쟁적 질문 & 케이블카, 에스컬레이터처럼 정해진 방향과 거리로만 움직이는 이동과, 자유롭게 방향을 바꿀 수 있는 이동 --- 우리 삶에는 어느 쪽이 더 많을까? \\
\end{tabularx}

\section{수업 흐름}
\begin{tabularx}{\linewidth}{@{}>{\centering\arraybackslash}p{1.6cm}X>{\raggedright\arraybackslash\small}p{3.1cm}@{}}
\toprule
\textbf{단계} & \textbf{활동 \& 발문} & \textbf{ATL / VTR} \\
\midrule
\textcolor{goldstrong}{\textbf{도입}}\newline\footnotesize 10분 &
\textbf{마음공부 (2분)} --- 유무념 대조.\newline
\textbf{생각 열기 (8분)} --- 장기판을 좌표평면으로 보고, 차(車)를 A에서 B로 두 수 만에 옮기는 방법을 말해 보게 한다(오른쪽 3칸, 위쪽 2칸). &
ATL: 사고(좌표 이동 직관)\newline VTR: See-Think-Wonder \\
\addlinespace
\textcolor{goldstrong}{\textbf{전개}}\newline\footnotesize 30분 &
\textbf{개념 형성 (12분)} --- 점 P$(x,y)\to$P'$(x+a,y+b)$ 확인 $\to$ 문제 1.\newline
\textbf{도형의 평행이동 (10분)} --- 도형 $F$ 위의 점을 이용해 $f(x-a,y-b)=0$을 유도(오개념: $f(x+a,y+b)=0$으로 착각하지 않도록 강조) $\to$ 예제 1, 문제 2.\newline
\textbf{실생활 적용 (8분)} --- 문제 3(캠핑장의 이동식 수영장) 짝 활동 --- 평행이동 전후 원의 방정식으로 이동 거리 $a,b$ 역산. &
ATT: 촉진자 역할\newline ATL: 사고·적용 \\
\addlinespace
\textcolor{goldstrong}{\textbf{정리}}\newline\footnotesize 10분 &
\textbf{개념 연결 질문 (5분)} --- ``점의 이동과 도형의 방정식의 이동, 부호가 반대인 이유를 내 말로 설명하면?''\newline
\textbf{성찰 (5분)} --- 감사 일기 1문장 + 다음 차시 예고(대칭이동). &
마음공부 · 감사 일기 \\
\bottomrule
\end{tabularx}

\begin{calloutbox}
\textbf{지도상의 유의점} --- 오개념 주의: $f(x,y)=0$을 $x$축 방향 $a$, $y$축 방향 $b$만큼 평행이동한 도형의 방정식은 $f(x-a,y-b)=0$이지, $f(x+a,y+b)=0$이 \textbf{아니다}. 도형 $F'$ 위의 점 P$'(x,y)$가 원래 도형 $F$ 위의 점 P$(x-a,y-b)$에서 왔다는 논리로 유도 과정을 반드시 짚어 준다.
\end{calloutbox}

\section{형성평가 \& 답안}
\begin{minipage}[t]{0.48\linewidth}
\begin{answerbox}
\textbf{문제 1} \quad $(-3,2)$만큼 평행이동\\
\small ⑴ $(-4,3)$ \quad ⑵ $(5,-1)$\\[4pt]
\textbf{\color{goldstrong}답} \; ⑴ $(-7,5)$ \quad ⑵ $(2,1)$
\end{answerbox}
\end{minipage}\hfill
\begin{minipage}[t]{0.48\linewidth}
\begin{answerbox}
\textbf{예제 1} \quad 원 $x^2+y^2=1$을 $(3,-1)$만큼 평행이동\\[4pt]
\textbf{\color{goldstrong}답} \; $(x-3)^2+(y+1)^2=1$
\end{answerbox}
\end{minipage}

\vspace{6pt}
\begin{answerbox}
\textbf{문제 2} \quad $(-2,4)$만큼 평행이동\\
\small ⑴ $x-4y-1=0$ \quad ⑵ $(x+1)^2+(y-5)^2=9$\\[4pt]
\textbf{\color{goldstrong}답} \; ⑴ $x-4y+17=0$ \quad ⑵ $(x+3)^2+(y-9)^2=9$
\end{answerbox}
\begin{answerbox}
\textbf{문제 3} \quad 원 $(x-6)^2+(y-4)^2=2$를 $(a,b)$만큼 평행이동하면 $x^2+y^2-2x-4y+3=0$\\
\textbf{\color{goldstrong}답} \; $a=-5$, $b=-2$
\end{answerbox}

\section{성취수준 (이 차시 범위)}
\begin{tabularx}{\linewidth}{@{}>{\bfseries\color{goldstrong}}p{0.9cm}X@{}}
\toprule
A & 도형의 평행이동을 탐구하고 방정식을 구하며, 실생활과 연결하여 문제를 해결할 수 있다. \\
B & 평행이동한 도형의 방정식을 구하고, 실생활과 연결할 수 있다. \\
C & 평행이동한 도형의 방정식을 구할 수 있다. \\
D & 평행이동한 도형의 방정식을 구할 수 있다(안내됨). \\
E & 평행이동한 점의 좌표를 구할 수 있다. \\
\bottomrule
\end{tabularx}

\section{다음 차시}
\begin{calloutbox}
\textbf{02 대칭이동 --- 1/2차시.} 점과 도형을 원점, $x$축, $y$축에 대하여 대칭이동하는 방법을 다룬다.
\end{calloutbox}
\end{document}
